\chapter{Why angle is not yet fitness}

\section{A geometric discrepancy has no inherent biological sign}

The quantities
\begin{equation}
d_C^2=\sum_j\sin^2\psi_j
\qquad\text{and}\qquad
I_V=1-\prod_j\cos^2\psi_j
\end{equation}
measure channel discrepancy. Neither contains a fitness scale, a habitat, a hybrid mean, or a hybrid covariance. Consequently neither can, by itself, determine whether hybrids are less fit, equally fit, or more fit than their parents.

\begin{cautionbox}
Calling \(I_V\) ``reproductive isolation'' without a fitness model would be circular. The channels would be declared incompatible because they overlap poorly, and poor overlap would then be treated as evidence of incompatibility. The missing bridge is a model of hybrid phenotypes evaluated on parental fitness surfaces.
\end{cautionbox}

\section{Hybrid phenotype distribution}

Let a hybrid phenotype be a random vector
\begin{equation}
\mathbf Z_H\sim(\muv_H,\Pp_H),
\end{equation}
where
\begin{equation}
\E[\mathbf Z_H]=\muv_H,
\qquad
\Cov(\mathbf Z_H)=\Pp_H.
\end{equation}
We do not initially require normality.

In environment \(e\), let the optimum and curvature be \(\thetav_e\) and \(\Hh_e\). Define the hybrid mean displacement
\begin{equation}
\mathbf d_e=\muv_H-\thetav_e.
\end{equation}
Write
\begin{equation}
\mathbf Z_H-\thetav_e
=\mathbf d_e+\boldsymbol\varepsilon,
\end{equation}
where
\begin{equation}
\E[\boldsymbol\varepsilon]=\Zero,
\qquad
\Cov(\boldsymbol\varepsilon)=\Pp_H.
\end{equation}

\section{Three different fitness summaries}

It is useful to distinguish three quantities.

\subsection{Individual log-fitness loss}
For one phenotype,
\begin{equation}
\ell_e(\mathbf z)
=-\log\frac{W_e(\mathbf z)}{W_{0e}}
=\half(\mathbf z-\thetav_e)^\top\Hh_e(\mathbf z-\thetav_e).
\end{equation}

\subsection{Expected log-fitness loss}
\begin{equation}
L^{\log}_{H,e}=\E[\ell_e(\mathbf Z_H)].
\end{equation}
This is algebraically simple and additive across independent log-fitness components.

\subsection{Arithmetic mean fitness}
\begin{equation}
\bar W_{H,e}=\E[W_e(\mathbf Z_H)].
\end{equation}
This is often the directly relevant demographic quantity, but its relationship to expected log fitness is nonlinear.

By Jensen's inequality,
\begin{equation}
-\log\E\left[\frac{W}{W_0}\right]
\le
\E\left[-\log\frac{W}{W_0}\right],
\end{equation}
because \(-\log x\) is convex. Thus the log of arithmetic mean fitness and expected log-fitness loss are not identical.

\begin{mathteach}
The geometric mean relative fitness is \(\exp[-L^{\log}]\). The arithmetic mean relative fitness is \(\bar W/W_0\). They agree only under special conditions or approximately when variation in log fitness is small.
\end{mathteach}

\section{Bounded isolation measures}

A log-scale barrier may be mapped to \([0,1)\) by
\begin{equation}
I_L=1-e^{-\Delta L^{\log}},
\end{equation}
where \(\Delta L^{\log}\) is hybrid log-load relative to a chosen parental benchmark. This is a geometric-mean interpretation.

An arithmetic-fitness isolation measure is
\begin{equation}
I_A=1-\frac{\bar W_H}{\bar W_{\mathrm{ref}}}.
\end{equation}
For weak load, both satisfy
\begin{equation}
I_L\approx\Delta L^{\log},
\qquad
I_A\approx-\log\left(\frac{\bar W_H}{\bar W_{\mathrm{ref}}}\right),
\end{equation}
but they should not be silently interchanged.
